\documentclass[12pt]{article}
\usepackage[a4paper, margin=2cm]{geometry}

\usepackage{fontspec}
\usepackage{polyglossia}
\setdefaultlanguage{russian}
\setmainlanguage{russian}
\setmainfont{Times New Roman}
\begin{document}

\begin{center}
\textbf{\LARGE ТЕХНИЧЕСКИЙ ОТЧЕТ}
\end{center}

\vspace{1cm}

% === ТЕКСТ ДОКУМЕНТА ВСТАВЛЯТЬ СЮДА ===
% Здесь должно быть описание проверяемого оборудования,
% проведенные измерения, анализ состояния и выводы.
% Допускается ссылка на нормативные документы и регламенты.

\textbf{Описание объекта проверки}

Объект проверки: промышленный узел технологической линии, предназначенный для перекачки рабочей жидкости. Модель: ПЗ-ЛТ, серия 4, заводской номер 000-00123. Назначение: контроль работоспособности узла и соответствия заявленным характеристикам в рамках плановой инспекции. 

Условия проведения испытаний: помещение цеха №3, температура 22 ± 2 °C, относительная влажность 50 ± 10%. Тестирование проводилось в период с 20.09.2026 по 22.09.2026.

Методика: визуальный осмотр, функциональные тесты, измерение следующих параметров: давление, температура, расход, вибрация, энергопотребление. Описание методики приведено в документе внутреннего контроля "Методика тестирования оборудования" (рег. № МТ-2024-05).

\textbf{Проведённые измерения}

Таблица 1. Результаты измерений (условные значения). Подтверждающие данные представлены в приложении.

\begin{tabular}{p{6cm} c p{3cm}}
Параметр & Единицы & Значение \\
Температура окружающей среды & °C & -- \\
Давление в системе & бар & -- \\
Уровень вибрации (X) & мм/с & -- \\
Энергопотребление & кВт & -- \\
Расход рабочей жидкости & л/мин & -- \\
Время тестирования & ч & -- \\
\end{tabular}

\textbf{Анализ состояния}

- Визуальный осмотр не обнаружил видимых дефектов на корпусе, вентилях и креплениях. Все соединения выглядят чистыми и надежно заизолированы.

- Функциональная проверка показала корректность выдачи управляющих сигналов и стабильность режимов работы в заданных пределах. Отклонения параметров от заданных порогов не превышали допустимых значений.

- Зафиксированы незначительные отклонения некоторых параметров на момент тестирования, которые требуют уточнения при повторной выборке данных и мониторинга.

\textbf{Выводы и рекомендации}

- Оборудование находится в рабочем состоянии и готово к дальнейшей эксплуатации при условии отсутствия динамических изменений в ближайшие 6 месяцев.

- Рекомендуется: 1) запланировать повторную проверку через 6 месяцев; 2) проверить крепления и электроподключение, а также состояние датчиков вибрации; 3) оформить акт по итогам инспекции и хранить в архиве технической документации.

\vfill

\noindent
Дата: \rule{5cm}{0.4pt} \hfill Подпись: \rule{5cm}{0.4pt}

\end{document}